\section{Global classification and recovery through the horizon}
\label{q3:global:section}

We continue with the finite-component class of
Theorem~\ref{q3:intro:rigidity}, with no prescribed boundary genera.
Starting from the complete static base and end estimates just proved,
we compactify one end of the double at its actual regularity and retain
all possible other ends while proving zero-mass rigidity. We then exclude the
interior null set.  This order preserves the information needed to
recover both the original metric and second fundamental form.

\begin{lemma}[The two conformal metrics]
\label{q3:static:conformal-completion}
Define
\[
h_\pm=\left(\frac{1\pm\lambda}{2}\right)^4h
\quad\text{on }\mathcal E.
\]
Both are scalar flat. The plus end satisfies
\(h_+=\delta+O_k(r^{-1-\epsilon})\), hence has zero ADM mass.
The minus end compactifies by one point to a nondegenerate
\(C^{1,\epsilon}\cap W^{2,p}_{\mathrm{loc}}\) metric for some
\(p>3\), with scalar curvature zero as a distribution at that point.
\end{lemma}

\begin{proof}
Scalar flatness follows from
Equation~\eqref{q3:static:static-equations} and the three-dimensional
conformal scalar-curvature formula, since \(1\pm\lambda\) are harmonic.
In terms of the preceding variables,
\begin{equation}
\label{q3:static:conformal-gamma}
h_+=\cosh^4(U/2)\gamma,\qquad
h_-=\sinh^4(U/2)\gamma.
\end{equation}
The first factor is \(1+O_k(r^{-2})\), proving the assertion for
\(h_+\). Its ADM flux is \(O(r^{-\epsilon})\), so its mass is zero.

For the minus end invert \(x=y/|y|^2\), write \(\rho=|y|\), and
put \(R(y)=I-2y\otimes y/\rho^2\).
The inversion differential is \(\rho^{-2}R(y)\), with
\(R(y)^TR(y)=I\).
Equation~\eqref{q3:static:gamma-expansion} gives
\[
I^*\gamma=\rho^{-4}
\bigl(I+O_k(\rho^{1+\epsilon})\bigr),
\qquad
U(I(y))=-M\rho+(b_1\cdot y)\rho
 +O_k(\rho^{2+\epsilon}).
\]
Derivatives of the angular matrix \(R(y)\) cost the corresponding
powers of \(\rho^{-1}\), so the displayed metric estimate retains
the symbol bounds. Dividing the expansion of \(\sinh(U/2)\) by
\(\rho\), and using \(\epsilon<1\), gives
\begin{equation}
\label{q3:static:inverted-minus}
I^*h_-=\left(\frac M2\right)^4
\left[
\left(1-\frac{4b_1\cdot y}{M}\right)I
 +O_k(\rho^{1+\epsilon})
\right].
\end{equation}
In particular, its only first-order term is linear and scalar.

For the remainder \(E(y)\), the bounds
\[
|E|\leq C\rho^{1+\epsilon},\quad
|DE|\leq C\rho^\epsilon,\quad
|D^2E|\leq C\rho^{\epsilon-1}
\]
show that \(E(0)=DE(0)=0\) extends it as \(C^{1,\epsilon}\).
To see the Hölder estimate between two points, use the gradient
bound when their separation is comparable to their radii; otherwise
integrate the Hessian along the short segment in a common annulus.
Moreover choose
\[
3<p<\frac3{1-\epsilon}.
\]
Then \(D^2E\in L^p\). Integration by parts across small coordinate
spheres has an error bounded by a constant times their area,
because the first derivatives are bounded.
It follows that the punctured classical derivatives are the weak
derivatives of the extension.
The Christoffel symbols are continuous with weak derivatives in
\(L^p\), so scalar curvature is its usual \(L^p\) expression
in \(\partial\Gamma\) and \(\Gamma\Gamma\).
It vanishes off the point and therefore also as a distribution on
the whole ball. There is no point curvature term.
\end{proof}

\subsection{Zero-mass rigidity with arbitrary remaining ends}
\label{q3:static:spin-subsection}

The static equations and end expansion are now established.  We have
not yet proved that the stationary field is timelike throughout the
original exterior.  The next step closes precisely that gap.
The completed conformal manifold can still have ends corresponding
to an interior zero set. We need a rigidity statement that permits
these ends, and also the regularity in
Lemma~\ref{q3:static:conformal-completion}.
The smooth arbitrary-end positive mass theorem is established in
\cite[Theorem B]{CecchiniZeidler2024}.
That formulation also credits the earlier spinor rigidity theorem of
Bartnik and Chru\'sciel \cite[Theorem~11.2]{BartnikChrusciel2005}.
In our scalar-flat situation the following direct spinor proof also
handles the compactified point. It uses Witten's
identity~\cite{Witten1981}; no asymptotic condition is imposed on any
other end.

\begin{lemma}[Scalar-flat rigidity at a faster-decaying end]
\label{q3:static:spin-rigidity}
Let \((W,\mathfrak g)\) be a connected orientable complete
three-dimensional Riemannian manifold without boundary.
Assume that its metric is locally \(C^{1,\alpha}\cap W^{2,p}\),
for some \(\alpha>0,p>3\), is scalar flat as a distribution,
and has an end on which it is smooth with
\[
\mathfrak g=\delta+O_2(r^{-1-\epsilon}),\qquad \epsilon>0.
\]
The other ends can be arbitrary.
Then \((W,\mathfrak g)\) is isometric to Euclidean space.
The asserted regularity is sufficient at a compactified point and
at a reflected hypersurface.
\end{lemma}

\begin{proof}
An orientable three-manifold is spin.
Choose a spin structure and use a smooth background metric to
identify its spinor bundle with that for \(\mathfrak g\).
The positive square-root change of orthonormal frames has the
stated metric regularity.
The spin connection consequently has continuous locally bounded
coefficients, with weak first derivatives in \(L^p_{\mathrm{loc}}\).
Write \(\nabla^S\) for that connection, \(D\) for the Dirac operator,
and \(c\) for Clifford multiplication.

The weak Lichnerowicz formula has the following integrated form for
every compactly supported \(H^1\) spinor \(\phi\):
\begin{equation}
\label{q3:static:lichnerowicz}
\|D\phi\|_2^2=\|\nabla^S\phi\|_2^2.
\end{equation}
Here and below norms and integrals are for \(\mathfrak g\).
To check the formula at the indicated regularity, first use a smooth
compactly supported spinor. Approximate the metric on its support
in \(C^1\) and \(W^{2,p}\) by smooth metrics, using the same
background-bundle identification.
The smooth Lichnerowicz formula passes to the limit: connection
coefficients converge uniformly and scalar curvatures converge
in \(L^p\).
The limiting scalar term is zero by the assumed distributional
scalar flatness and the \(L^p\) curvature expression.
Density then proves Equation~\eqref{q3:static:lichnerowicz} for compact
\(H^1\) spinors.
Thus no integration surface is retained around the compactified point.

We next establish the coercivity needed to complete compact spinors.
For a compactly supported scalar function \(f\) on the AE end,
one-dimensional integration by parts along coordinate rays gives
\begin{equation}
\label{q3:static:radial-hardy}
\int_{r_0}^\infty f^2\dd r
\leq4\int_{r_0}^\infty r^2|\partial_rf|^2\dd r.
\end{equation}
The inner boundary term has the favorable sign.
Apply this to \(f=|\phi|\) and integrate over angles.
Metric comparability and
\(|\dd|\phi||\leq|\nabla^S\phi|\) give
\begin{equation}
\label{q3:static:spin-hardy}
\int_{\mathrm{AE}}r^{-2}|\phi|^2\dd V_{\mathfrak g}
\leq C\|\nabla^S\phi\|_2^2.
\end{equation}
For any fixed compact \(K\subset W\), connect it to a fixed end
annulus by a relatively compact connected domain.
Poincaré's inequality with the norm on that annulus retained gives
\[
\||\phi|\|_{L^2(K)}
\leq C_K\left(
\|\dd|\phi|\|_{L^2(\text{domain})}
 +\|\phi\|_{L^2(\text{annulus})}\right).
\]
For example, this version follows from the ordinary Poincaré
inequality after subtracting the mean; the observed annulus
controls that mean because it has positive measure.
Equation~\eqref{q3:static:spin-hardy} controls its last term.
We have proved
\begin{equation}
\label{q3:static:local-coercivity}
\|\phi\|_{L^2(K)}\leq C_K\|\nabla^S\phi\|_2.
\end{equation}
Connectedness is the only global input in this propagation.

Let \(\mathcal H\) be the completion of compactly supported smooth
spinors in the norm \(\|\nabla^S\phi\|_2\).
Equation~\eqref{q3:static:local-coercivity} embeds \(\mathcal H\)
in \(L^2_{\mathrm{loc}}\), and the locally bounded connection
then embeds it in \(H^1_{\mathrm{loc}}\).
Both \(\nabla^S:\mathcal H\to L^2\) and
\(D:\mathcal H\to L^2\) are well-defined, and
Equation~\eqref{q3:static:lichnerowicz} makes the latter an isometry.
Its range is therefore closed.
Equation~\eqref{q3:static:spin-hardy} also passes to this completion.
The local constants may depend on the chosen compact set; no uniform
estimate over the remaining ends is required.

Fix the ordinary orthonormalized coordinate frame on the distinguished
end and a smooth cutoff \(\theta\), equal to one near infinity and
zero before a slightly smaller end neighborhood.
For each constant spinor \(z\in\mathbb C^2\), set
\(\psi_0=\theta z\).
Because the connection there is \(O(r^{-2-\epsilon})\),
\(\nabla^S\psi_0,D\psi_0\in L^2\).
Orthogonally project \(D\psi_0\) onto the closed subspace
\(D\mathcal H\) and choose \(\eta\in\mathcal H\) such that
\[
D\eta=-\operatorname{proj}_{D\mathcal H}(D\psi_0).
\]
Set \(\psi=\psi_0+\eta\), \(w=D\psi\).
Then \(w\in L^2\) and
\[
\int\langle w,D\phi\rangle\dd V_{\mathfrak g}=0
\quad\text{for every compact smooth }\phi.
\]
Thus \(Dw=0\) weakly.
Local elliptic regularity for a first-order elliptic operator with
locally Lipschitz principal coefficients and bounded lower-order
coefficients gives \(w\in H^1_{\mathrm{loc}}\).
This regularity can also be seen by freezing the principal
coefficients on small coordinate balls: the constant-coefficient
Dirac estimate controls one derivative in \(L^2\), the oscillation
is absorbed, and commutators of mollification are bounded by the
local Lipschitz norm. Apply the estimate to regularized \(w\)
and pass to a weak limit.
The metric in the lemma has more than this coefficient regularity.

Completeness supplies compactly supported Lipschitz distance cutoffs
\(\chi_R\), equal to one on the radius-\(R\) ball and with
\(|\dd\chi_R|\leq C/R\).
Indeed a complete locally compact Riemannian length space is proper,
and a piecewise linear cutoff of its distance has these properties.
Equation~\eqref{q3:static:lichnerowicz} applies to \(\chi_Rw\), and
the weak equation gives
\[
\|\nabla^S(\chi_Rw)\|_2^2
=\|D(\chi_Rw)\|_2^2
=\|c(\dd\chi_R)w\|_2^2
\leq \frac C{R^2}\|w\|_2^2.
\]
On every fixed compact set the cutoff eventually equals one.
It follows that \(\nabla^S w=0\).
A parallel spinor has constant length on the connected manifold;
the distinguished end has infinite volume, so \(w\in L^2\)
forces \(w=0\).
No condition at another end entered this argument.

It remains to show that \(\psi\), rather than just \(D\psi\), is
parallel. Define the continuous sesquilinear form
\[
\mathcal B(a,b)=
\int\left(\langle\nabla^Sa,\nabla^Sb\rangle
-\langle Da,Db\rangle\right)\dd V_{\mathfrak g}
\]
on fields for which the indicated derivatives are in \(L^2\).
For compact \(\phi\), the polarized Lichnerowicz identity,
with one slot compactly supported, gives
\(\mathcal B(\psi_0,\phi)=0\).
Approximate \(\eta\) in \(\mathcal H\) by compact spinors.
These limits use the global derivative norms; no boundary flux of
\(\eta\) at another end is invoked.
Continuity and Equation~\eqref{q3:static:lichnerowicz} imply
\[
\mathcal B(\psi_0,\eta)=0,\qquad
\mathcal B(\eta,\eta)=0,\qquad
\mathcal B(\psi,\psi)=\mathcal B(\psi_0,\psi_0).
\]
The last expression is zero. To see this directly, integrate the
Lichnerowicz identity for \(\psi_0\) up to a large end sphere.
There is no contribution at another end because \(\psi_0\)
vanishes there. On that sphere \(\psi_0=O(1)\) and
\(\nabla^S\psi_0,D\psi_0=O(r^{-2-\epsilon})\);
the boundary integral is \(O(r^{-\epsilon})\) and tends to zero.
This is also the zero ADM boundary term of the spin proof.
Since \(D\psi=w=0\), we conclude that
\(\|\nabla^S\psi\|_2^2=0\).

The parallel equation gives each resulting spinor a
\(C^{1,\beta}_{\mathrm{loc}}\) representative for some \(\beta>0\).
Indeed, \(H^1\subset L^6\) and the locally bounded connection first
give \(W^{1,6}\), hence local H\"older continuity; the locally H\"older
connection coefficients then make its first derivatives H\"older.
Write \(\psi_z\) for this parallel spinor. The fixed cutoff,
orthogonal projection, and inverse of \(D:\mathcal H\to D\mathcal H\)
make \(z\mapsto\psi_z\) complex linear.
For any point \(p\in W\), consider
\[
 E_p:\mathbb C^2\longrightarrow\mathbb S_p,
 \qquad E_p(z)=\psi_z(p),
\]
where \(\mathbb S_p\) is the complex spinor fiber.
If \(E_p(z)=0\), parallelism and connectedness give \(\psi_z=0\)
everywhere. Then \(\eta=-\theta z\), which for \(z\ne0\) has infinite
\(r^{-2}\)-weighted \(L^2\) norm on the distinguished end,
contrary to Equation~\eqref{q3:static:spin-hardy}.
Thus \(E_p\) is injective for every \(p\), and it is an isomorphism
because the complex spinor rank is two. The images of a basis of
\(\mathbb C^2\) form a global parallel spinor frame.

Its Hermitian Gram matrix is constant and positive definite;
a constant change of basis makes the frame orthonormal.
The weak curvature equations annihilate this frame, so the spin
curvature vanishes. Faithfulness of the spin representation on
\(\mathfrak{spin}(3)\) gives zero Riemann curvature.
Moreover, Clifford multiplication identifies the real tangent bundle
with the traceless skew-Hermitian endomorphisms of the spinor bundle.
Three fixed orthonormal matrices in this parallel spinor frame,
with matrix inner product \(-\tfrac12\operatorname{tr}(AB)\),
therefore give a global parallel orthonormal tangent frame.
The dual coframe is closed by torsion-freeness and integrates
locally to coordinates in which the metric is exactly Euclidean.
This also removes any apparent metric singularity at the
compactified point.

Finally, completeness makes the universal flat cover Euclidean
\(\R^3\). The global parallel frame makes its deck transformations
translations. A nontrivial discrete translation group has rank
at least one, and a quotient by it has volume growth at most
quadratic. The distinguished AE end gives a cubic lower bound:
radial paths in that end put a Euclidean annulus of radius
comparable to \(R\) inside a metric ball of radius \(CR\).
The two volume bounds are incompatible.
The deck group is therefore trivial, proving the lemma.
\end{proof}

\begin{lemma}[The infinity component has no interior missing boundary]
\label{q3:static:no-internal-null}
The component \(\mathcal E\) is all of \(\Omega\).
In particular \(N>0\) everywhere in the open exterior and the
stationary metric is vacuum there.
\end{lemma}

\begin{proof}
The complete base of Lemma~\ref{q3:static:base-completion} reaches every
component of $S$. Take its complete two-copy double, use $h_+$ on
the plus copy and $h_-$ on the minus copy, and identify corresponding
boundary points. The signed reflected lapse expresses the joined
metric as
\[
q=\left(\frac{1+\lambda_{\mathrm{signed}}}{2}\right)^4h_{\mathrm d}.
\]
It is smooth at each positive-lapse join and $C^2$ at each zero-lapse
join. The matched second jets make its scalar curvature zero across
every join, as well as on the two open sides. Compactify the unique
AF end of the minus copy by
Lemma~\ref{q3:static:conformal-completion}.

We verify completeness of this conformal manifold without placing any
restriction on the number or regularity of the possible internal-null
ends. On the plus copy the factor $(1+\lambda)/2$ is at least $1/2$.
Fix a sufficiently large original AF tail. On its complement the
continuous function $\sqrt N$ is defined on a compact subset of the
original closed exterior, equals zero on $S$ and the internal zero
set, and is strictly less than one elsewhere by
Lemma~\ref{q3:static:base-completion}. Its maximum on this compact set
is therefore some $a_0<1$. Thus on the entire corresponding minus
remainder $(1-\lambda)/2\geq(1-a_0)/2>0$ uniformly. On the double
away from the minus AF tail the conformal metric is consequently
bounded below by a fixed positive multiple of the complete doubled
base metric. All joins are already attached. A finite-length escape
can therefore only run down the minus AF tail, and its proved
nondegenerate one-point compactification supplies its limit. More
explicitly, truncate that tail at a large sphere. The other side of
the sphere is complete with compact boundary by the uniform lower
bound, and the compactified tail is a compact metric region with
that same boundary; their gluing is complete.

The conformal manifold is connected, orientable, and without boundary.
It has local $C^{1,\alpha}\cap W^{2,p}$ regularity for some
$\alpha>0$ and $p>3$: the finitely many joins are at least $C^2$,
and the sole compactification point has the regularity established
in Lemma~\ref{q3:static:conformal-completion}. Its scalar curvature is
zero distributionally, including at that point, and its distinguished
plus end has the fast decay required by
Lemma~\ref{q3:static:spin-rigidity}. That lemma makes this entire
manifold Euclidean space.

Choose a sufficiently large sphere in the distinguished plus end.
Under the Euclidean isometry it is a compact embedded sphere. Its
end tail is the unbounded component of the complement: paths from
large end radii to that sphere have lengths tending to infinity,
and the tail has precisely that sphere as frontier. The other
component is bounded and has compact closure. A sequence in
$\mathcal E$ approaching an interior point with $N=0$ lies outside
the chosen end tail and escapes every compact set of the complete
conformal manifold. Indeed it cannot converge at an ordinary plus
point by continuity of $N$, cannot converge on any joined component
because every original boundary collar has no interior zero, and
cannot converge at a point in the open minus copy or its compactified
end. Such a sequence is impossible in the compact Euclidean
complement. Thus $\mathcal E$ has no interior boundary in $\Omega$.

It is a nonempty open and closed subset of connected $\Omega$, and
hence equals $\Omega$. Vacuum now follows from
Proposition~\ref{q3:variation:adjoint}, since $N>0$ everywhere in
$\Omega$.
\end{proof}

\begin{lemma}[Equality forces a single spherical boundary component]
\label{q3:static:one-boundary-component}
In the Euclidean conformal double, every original component $S_i$
is a round sphere of radius $a=(8\kappa)^{-1}$ and has original area
$|S_i|_g=16\pi m^2=|S|_g$. Consequently $\ell=1$.
\end{lemma}

\begin{proof}
The completed conformal double in the proof of
Lemma~\ref{q3:static:no-internal-null} is Euclidean. In particular its
plus metric $q=h_+$ is flat up to every original boundary component.
Put $\phi=(1+\lambda)/2$. On each $S_i$,
Lemma~\ref{q3:static:base-completion} gives
\[
\phi=\tfrac12,\qquad
\partial_{\nu_h}\log\phi=\kappa,\qquad
\mathrm{II}_h=0,\qquad h|_{TS_i}=g|_{TS_i}.
\]
The conformal second-fundamental-form formula, with normal pointing
into the plus side, is therefore
\[
\mathrm{II}_q
 =\phi^2\bigl(\mathrm{II}_h+
 2\partial_{\nu_h}\log\phi\,h|_{TS_i}\bigr)
 =8\kappa\,q|_{TS_i},
 \qquad q|_{TS_i}=\tfrac1{16}g|_{TS_i}.
\]
Let $F_i$ be the Euclidean position vector on $S_i$ and $\nu_i$
its unit normal into the plus side. Its shape operator is
$8\kappa I$, so the tangential differential of $F_i-a\nu_i$ is
zero for $a=(8\kappa)^{-1}$. Connectedness makes $F_i-a\nu_i$
a constant point $c_i$. Thus the image lies on the sphere of radius
$a$ centered at $c_i$, with $\nu_i$ its outward radial normal.
The image is open in that sphere by local invertibility and closed
by compactness, hence is the whole sphere. The Euclidean embedding
therefore identifies $S_i$ diffeomorphically with it.
Since $q|_{TS_i}=g|_{TS_i}/16$, its area satisfies
\begin{equation}
\label{q3:static:individual-boundary-area}
|S_i|_g=16|S_i|_q=64\pi a^2
 =\frac{\pi}{\kappa^2}=16\pi m^2=|S|_g.
\end{equation}
Summing over the finite family yields
$|S|_g=\sum_{i=1}^{\ell}|S_i|_g=\ell|S|_g$. The boundary is nonempty
and $|S|_g>0$, hence $\ell=1$. In particular no connectedness or spherical
topology assumption was needed before this step.
\end{proof}

\begin{lemma}[Identification of the base and its mass]
\label{q3:static:schwarzschild-base}
The regular completion of \((\Omega,h,\lambda)\) is the
Schwarzschild static exterior
\begin{equation}
\label{q3:static:isotropic-base}
h=\left(1+\frac a\rho\right)^4\delta,\qquad
\lambda=\frac{1-a/\rho}{1+a/\rho},\qquad
\rho\geq a,
\end{equation}
where \(a=m/2\).
The completion is homeomorphic to the original \(\Mext\),
and \(S\) is identified diffeomorphically with its horizon sphere
in the regular base boundary structure.
\end{lemma}

\begin{proof}
Lemma~\ref{q3:static:one-boundary-component} identifies the boundary
in the Euclidean conformal double as one round sphere of radius
$a=(8\kappa)^{-1}$. The plus side is its outside: its normal is the
outward radial normal, its entire boundary is this sphere, and it
contains the distinguished unbounded end. Center Euclidean
coordinates at this sphere and write $\rho$ for their radius.

In the resulting Euclidean coordinates, write
\(\psi=2/(1+\lambda)\), so \(h=\psi^4\delta\).
The scalar-flat conformal equation gives \(\Delta_\delta\psi=0\)
outside the sphere, with \(\psi=2\) on it and \(\psi\to1\)
at infinity. The function \(1+a/\rho\) has the same data.
Their difference vanishes by the exterior maximum principle.
This proves Equation~\eqref{q3:static:isotropic-base}.
Its Schwarzschild mass is $2a=(4\kappa)^{-1}=m$, since
$\kappa=1/(4m)$. In particular $a=m/2$.

The regular base collar has the same underlying boundary topology
as the original collar, even when its coordinate is
\(\lambda=\sqrt N\) instead of the original defining function \(N\).
Together with the unchanged interior this identifies its completion
homeomorphically with \(\Mext\).
An interior base isometry extends uniquely to the metric completions.
On the regular boundary structures it is smooth: its boundary
restriction preserves induced lengths, hence is a smooth boundary
isometry, and normal geodesic coordinates extend it smoothly to
a collar.
Thus the completion has the claimed global topology, and
\(\Omega\cong(a,\infty)\times S^2\) is simply connected.
\end{proof}

\subsection{The time graph and the Kruskal attachment}
\label{q3:static:embedding-subsection}

\begin{proof}[Proof of Proposition~\ref{q3:static:classification}]
Lemma~\ref{q3:static:one-boundary-component} proves $\ell=1$ and the
spherical topology of $S$. Lemma~\ref{q3:static:no-internal-null} proves
\(N>0\) on \(\Omega\), and
Lemma~\ref{q3:static:schwarzschild-base} gives the global static
base with mass \(m\).
For the reconstruction write \(M=m\), so \(\kappa=1/(4M)\).
Simple connectivity and Lemma~\ref{q3:static:base-completion}
give a global smooth function \(H\) with
\begin{equation}
\label{q3:static:time-primitive}
\dd H=-A.
\end{equation}
In static Schwarzschild coordinates use time \(T=t+H(x)\).
Equations~\eqref{q3:static:base-definitions} and
\eqref{q3:static:time-primitive} give the exact identity
\begin{equation}
\label{q3:static:reconstructed-metric}
-N\dd T^2+h
=-u^2\dd t^2+g_{ij}(\dd x^i+X^i\dd t)
                         (\dd x^j+X^j\dd t).
\end{equation}
The graph \(T=H(x)\), corresponding to \(t=0\), therefore induces
the original metric \(g\).
The stationary future unit normal is \(n=(\partial_t-X)/u\).
With the second-fundamental-form convention of the theorem,
\begin{equation}
\label{q3:static:reconstructed-k}
K_{\mathrm{graph}}
=\frac1{2u}\bigl(\partial_tg-\mathcal L_Xg\bigr)
=-\frac1{2u}\mathcal L_Xg
=K.
\end{equation}
This calculation recovers the original \(K\), including a
non-time-symmetric slice.

We next establish smoothness in the original boundary structure.
If \(u|_S=0\), Equation~\eqref{q3:static:double-norm} gives
\[
A=\frac{Y^\flat}{a^2-s^2|Y|_g^2}.
\]
This is smooth in the original collar.
The primitive \(H\) therefore extends smoothly and finitely to \(S\).
For instance, fix its values on one interior collar section and
integrate its smooth normal derivative to the boundary.
Tangential derivatives of this extension agree with those prescribed
by \(\dd H=-A\), by continuity from the interior.

If \(u|_S>0\), the numerator
\(X^\flat-(2\kappa)^{-1}\dd N\) vanishes as a one-form on \(S\),
by Equation~\eqref{q3:static:simple-norm}.
Smooth division by the defining function \(N\) gives
\begin{equation}
\label{q3:static:log-time}
A=\frac1{2\kappa}\dd\log N+\beta,\qquad
H=-\frac1{2\kappa}\log N+B,
\end{equation}
where \(\beta\) and \(B\) are smooth in the original collar.
The smoothness of \(B\) follows by applying the same normal-integration
argument to its differential \(-\beta\).
In this case \(H\to+\infty\) at the boundary, so static time itself
is not a regular attachment coordinate.

Let \(r\) be the Schwarzschild area radius. On the whole exterior
\[
N=1-\frac{2M}{r},\qquad
r=\frac{2M}{1-N}.
\]
Normalize the tortoise coordinate and Kruskal coordinates by
\begin{equation}
\label{q3:static:kruskal-definitions}
\begin{aligned}
r_*&=r+2M\log\left(\frac r{2M}-1\right),\\
\mathsf U&=-e^{-\kappa(T-r_*)},&
\mathsf V&= e^{\kappa(T+r_*)}.
\end{aligned}
\end{equation}
On the right exterior \(\mathsf U<0,\mathsf V>0\), and
the Schwarzschild metric extends as
\[
\mathbf g_M
=-\frac{32M^3}{r}e^{-r/(2M)}
   \dd\mathsf U\,\dd\mathsf V+r^2\dd\omega^2.
\]
An exact useful form of Equation~\eqref{q3:static:kruskal-definitions}
is
\begin{equation}
\label{q3:static:kruskal-factor}
\mathsf U=-Q(N)\sqrt N\,e^{-\kappa T},\qquad
\mathsf V= Q(N)\sqrt N\,e^{\kappa T},\qquad
Q(N)=\frac{\exp\!\bigl(1/[2(1-N)]\bigr)}{\sqrt{1-N}}.
\end{equation}
The function \(Q\) is smooth and positive near zero, with
\(Q(0)=\sqrt e\).

If \(u|_S=0\), both \(H\) and
\(\sqrt N=s\sqrt{a^2-s^2|Y|^2}\) are smooth.
Substitution into Equation~\eqref{q3:static:kruskal-factor} gives
\(\mathsf U=\mathsf V=0\) on \(S\) and
\[
\partial_s\mathsf U|_S
=-\sqrt e\,\kappa e^{-\kappa H|_S}<0,\qquad
\partial_s\mathsf V|_S
=\sqrt e\,\kappa e^{\kappa H|_S}>0.
\]
This is a smooth attachment at the bifurcation sphere.
If \(u|_S>0\), Equation~\eqref{q3:static:log-time} instead gives
\begin{equation}
\label{q3:static:future-attachment}
\mathsf U=-Q(N)N e^{-\kappa B},\qquad
\mathsf V=Q(N)e^{\kappa B}.
\end{equation}
These are smooth in the original collar;
\(\mathsf U\) vanishes simply and
\(\mathsf V|_S=\sqrt e\,e^{\kappa B|_S}>0\).
Thus the attachment is on the future horizon, with each generator
met once after the angular identification.

The angular regularity requires a further check when
\(u|_S>0\), because the regular base coordinate was \(\lambda\),
not the original \(N\).
On the smooth reflected collar in
Equation~\eqref{q3:static:even-base}, normal projection onto its fixed
boundary is invariant under \(\lambda\mapsto-\lambda\).
This follows from uniqueness of the normal geodesics and reflection
invariance of the metric.
Schwarzschild angular coordinates are constant along these normal
geodesics, and their boundary values are the smooth round-sphere
identification from Lemma~\ref{q3:static:schwarzschild-base}.
They are therefore smooth even functions of \(\lambda\).
A smooth even function, with smooth tangential parameters, is a
smooth function of \(N=\lambda^2\) for \(N\geq0\):
all odd Taylor coefficients vanish, and repeated application of
\((2\lambda)^{-1}\partial_\lambda\) to the Taylor remainder gives
continuous derivatives of every order at zero.
This proves the required angular smoothness in the original collar.
In the zero boundary-lapse case the original collar is already a
smooth base collar, so its completed base isometry supplies the
angular extension directly.

We have constructed a smooth map \(\iota:\Mext\to\) Schwarzschild.
Equation~\eqref{q3:static:reconstructed-metric} extends to \(S\)
by continuity, hence \(\iota^*\mathbf g_M=g\) there as well.
The positive definiteness of \(g\) proves that \(\dd\iota\) is
injective at every boundary point.
Interior injectivity follows from the global base isometry;
boundary injectivity follows from the sphere identification.
The boundary image, at \(r=2M\), is disjoint from the open image,
where \(r>2M\).
The map is proper: a compact subset of Schwarzschild has bounded
area radius, and its inverse image is closed in the compact base
region corresponding to
\([2M,R]\times S^2\).
Thus the injective immersion is an embedding.

The future normal extends smoothly even where \(u=0\).
To avoid taking a limit of the expression involving \(u^{-1}\),
choose the smooth future timelike Kruskal vector
\(Z=\partial_{\mathsf U}+\partial_{\mathsf V}\) near the boundary.
Let \(Z^\top\) be its tangential projection along the embedded
spacelike hypersurface and set
\[
\widehat n=
\frac{Z-Z^\top}
{\sqrt{-\mathbf g_M(Z-Z^\top,Z-Z^\top)}}.
\]
The denominator is positive, and this is a smooth future unit
normal. On the open exterior it agrees with the normal used in
Equation~\eqref{q3:static:reconstructed-k}.
Its second fundamental form is smooth and equals \(K\) in the
interior, so equality extends to \(S\).

There is also a smooth spacelike extension of the image through
its boundary. In the positive boundary-lapse case, use
\((\mathsf U,\omega)\) as hypersurface coordinates; its normal
derivative is nonzero by
Equation~\eqref{q3:static:future-attachment}.
The remaining null coordinate is a smooth graph function and can
be extended across \(\mathsf U=0\).
In the zero boundary-lapse case use
\((\chi,\omega)\), where
\(\chi=(\mathsf V-\mathsf U)/2\); the displayed normal derivatives
give \(\partial_s\chi>0\).
Again extend the remaining graph function across zero.
Smooth one-sided graph functions extend across a smooth boundary,
and compactness of \(S\), together with openness of the spacelike
condition, preserves spacelikeness on a sufficiently short collar.
This extension makes no assertion about the original data behind
\(S\).

Finally, the exact inverse-metric identity gives
\begin{equation}
\label{q3:static:end-slope}
|\dd H|_h^2
=\left|\frac{X^\flat}{N}\right|_C^2
=\frac{|X|_g^2}{u^2N}
\longrightarrow\frac{|b|^2}{(b^0)^2}<1.
\end{equation}
Choose a constant \(a_0<1\) larger than the limiting slope norm.
Along radial rays in the Schwarzschild base, uniformly in angle,
\[
|H(r,\omega)|
\leq C+a_0\int_R^r
\frac{\dd\rho}{\sqrt{1-2M/\rho}}
=a_0r+O(\log r).
\]
Since \(r_*=r+O(\log r)\), it follows that
\[
H-r_*\longrightarrow-\infty,\qquad
H+r_*\longrightarrow+\infty.
\]
This is approach to the chosen Schwarzschild spatial infinity.
It allows a nonzero asymptotic tilt and does not impose \(P=0\)
or \(K=0\).
All conclusions of Proposition~\ref{q3:static:classification} follow.
\end{proof}

Propositions~\ref{q3:variation:adjoint} and~\ref{q3:static:classification}
prove Theorem~\ref{q3:intro:rigidity}. They also prove
Theorem~\ref{q3:intro:connected-rigidity}: the connected-case argument
in Lemma~\ref{q3:variation:area-support} supplies the same localized
multiplier identity, and all subsequent adjoint, static, and
reconstruction steps use that identity rather than
partial-coincidence outermostness.
